\subsection{A bound for one five-site profile}
\label{sec:five-site-profile}

The general five-site constant is not determined here.  A smaller bound
is available when the five values of the input $g$ have the prescribed
proportions $9:5:5:5:5$.  These are values of $g$, rather than values of
the separate kernel $k$.

\begin{theorem}\label{thm:five-profile-resolvent}
Let $H=\{h_0,h_1,h_2,h_3,h_4\}\subset\mathbb Z$ have five distinct
elements.  Suppose $f,g$ are nonzero, nonnegative and finitely supported,
and for some $t>0$,
\[
 g(h_0)=t,\qquad g(h_i)=\frac59t\quad(1\le i\le4).
\]
Then
\begin{equation}\label{eq:five-profile-bound}
 \frac{\|f\star\widetilde{g1_H}\|_1}{\|f\star g\|_1}
 \le \gamma,\qquad
 \gamma=\frac{3143504}{1566459}.
\end{equation}
The positions $h_i$ and the values of $f$ are unrestricted.
\end{theorem}

\begin{proof}
Restricting $g$ to $H$ leaves the numerator unchanged and can only
decrease the denominator.  Scaling $g$ and translating its support
also leave the ratio unchanged.  We may therefore take $h_0=0$,
$g_0=1$, $g_i=r=5/9$ for $1\le i\le4$, and $g=0$ outside $H$.
Put
\[
 G=\sum_{i=0}^4g_i=\frac{29}{9},\qquad
 P=\sum_x\max_{0\le i\le4}g_i f(x-h_i),\qquad
 Q=\sum_x\max_{0\le i\le4}g_i f(x+h_i).
\]
Every sum over $x$ is over $\mathbb Z$.

\medskip\noindent\textbf{Translation defects.}
Let $\mathcal S$ consist of the four three-element subsets of
$\{1,2,3,4\}$.  Define
\[
 \begin{aligned}
 M(x)&=\max\bigl(f(x),r f(x-h_1),\ldots,r f(x-h_4)\bigr),\\
 E_S(x)&=\max\bigl(f(x),\max_{i\in S}r f(x-h_i)\bigr),&
 P_S&=\|E_S\|_1,\\
 d_S(x)&=\bigl(f(x)-\max_{i\in S}r f(x-h_i)\bigr)_+,&
 D_S&=\|d_S\|_1.
 \end{aligned}
\]
With $\delta=(1-r)^3=64/729$, set
\[
 R_S=D_S-\delta P_S,\qquad
 R=\sum_{S\in\mathcal S}R_S,\qquad
 Z=\sum_{S\in\mathcal S}(P-P_S).
\]
We next prove $R_S\ge0$ and an estimate whose error depends on $R$
and $Z$.

For $S=\{i,j,\ell\}$ consider the nonnegative series
\[
 F_S(x)=
 \sum_{a,b,c\ge0}r^{a+b+c}
 d_S(x-a h_i-b h_j-c h_\ell).
\]
This auxiliary function need not be finitely supported.  Tonelli's
theorem gives $\|F_S\|_1=(1-r)^{-3}D_S<\infty$.
The inequality
\[
 f(x)\le d_S(x)+\max_{i\in S}r f(x-h_i)
\]
can be iterated along paths of shifts in $S$.  At depth $n$, the
remaining term is at most $r^n\|f\|_\infty$.  Along any one path,
the nonnegative defect terms use distinct multi-indices and are
bounded by the series defining $F_S$.  Letting $n$ tend to infinity
therefore gives $f\le F_S$.  Shifting a summation index also gives
\[
 rF_S(x-h_j)\le F_S(x)\qquad(j\in S).
\]
It follows that $E_S\le F_S$, and hence
\begin{equation}\label{eq:five-profile-resolvent-excess}
 R_S\ge0,\qquad
 \|F_S-E_S\|_1=\delta^{-1}R_S.
\end{equation}

Write $u_S=M-E_S\ge0$.  For $j\in S$, the preceding inequalities imply
\[
 \bigl(rM(x-h_j)-M(x)\bigr)_+
 \le r u_S(x-h_j)+F_S(x)-E_S(x).
\]
Choose distinct triples $S_j$ containing $j$: for example, omit
$j+1$ cyclically from $\{1,2,3,4\}$.  Each triple then occurs once.
Thus, with
\[
 e_j(M)=\sum_x\bigl(rM(x-h_j)-M(x)\bigr)_+,
\]
we have
\begin{equation}\label{eq:five-profile-translation-error}
 \sum_{j=1}^4e_j(M)\le rZ+\delta^{-1}R.
\end{equation}

For any nonnegative summable $u$, put
$b_j(x)=(r u(x-h_j)-u(x))_+$ and $e_j(u)=\sum_x b_j(x)$.
The definition gives
\[
 \min\{u(x),r u(x+h_j)\}
 \ge r^2u(x)-r b_j(x+h_j).
\]
Subtract this minimum from $r u(x+h_j)$ and sum over $x$.
Bounding a maximum by its initial coordinate plus the four positive
increments then gives
\[
 \sum_x\max\bigl(u(x),r u(x+h_1),\ldots,r u(x+h_4)\bigr)
 \le \alpha\|u\|_1+r\sum_{j=1}^4e_j(u),
 \qquad
 \alpha=1+4r(1-r)=\frac{161}{81}.
\]
Apply this to $u=M$, using $f\le M$, $\|M\|_1=P$ and
\eqref{eq:five-profile-translation-error}.  We obtain
\begin{equation}\label{eq:five-profile-stability}
 Q\le\alpha P+\tau R+r^2Z,\qquad
 \tau=r\delta^{-1}=\frac{405}{64}.
\end{equation}

\medskip\noindent\textbf{A tree estimate retaining the same defects.}
For $v=(v_0,\ldots,v_4)\ge0$, define
\[
 \Phi_{ij}(v)=\min(v_i,v_j)
 +\min\left(\frac{g_j}{g_i}v_i,\frac{g_i}{g_j}v_j\right),
 \qquad
 \kappa_{ij}=\frac{\min(g_i,g_j)}{g_i+g_j}.
\]
We have
\[
 \kappa_{ij}\Phi_{ij}(v)\le\min(v_i,v_j).
\]
For any spanning tree on the five indices, root the tree at a
coordinate attaining $\max_i v_i$ and charge each edge to its child.
Every other coordinate is charged once.  Therefore
\[
 \max_i v_i+\sum_{\{i,j\}\in T}\kappa_{ij}\Phi_{ij}(v)
 \le\sum_i v_i.
\]
Take a mixture of stars whose center probabilities are
\[
 p_0=\frac{103}{141}=1-2\beta,\qquad
 p_i=\frac{19}{282}=\frac{\beta}{2}\quad(1\le i\le4),
 \qquad\beta=\frac{19}{141}.
\]
Let
\[
 w_{ij}=(p_i+p_j)\kappa_{ij},\qquad
 W(v)=\sum_{i<j}w_{ij}\Phi_{ij}(v),\qquad
 q_-(v)=\sum_i v_i-W(v).
\]
The tree inequality gives $\max_i v_i\le q_-(v)$.
For the upper estimate use
\[
 a=1-\frac32\beta=\frac{75}{94},\qquad
 B=ar+2\beta=\frac{67}{94},\qquad
 A=4ar-\beta=\frac{77}{47},
\]
\[
 q_+(v)=Av_0+B\sum_{i=1}^4v_i-W(v).
\]
The root--leaf edge marginals are $a$ and the leaf--leaf marginals
are $\beta$.  The function $q_+$ is convex and positively
homogeneous, so
$q_+(v)\le\int_0^\infty q_+(1_{\{v_i\ge s\}})\,ds$.

For each triple $S$, write
\[
 \Delta_S(v)=\bigl(v_0-\max_{i\in S}v_i\bigr)_+,\qquad
 M_S(v)=\max\bigl(v_0,\max_{i\in S}v_i\bigr),
\]
and set
\[
 z(v)=4\max_i v_i-\sum_{S\in\mathcal S}M_S(v).
\]
The quantities $\max_i v_i$, $M_S(v)$, $\Delta_S(v)$ and $z(v)$
equal the integrals of their values on these indicator layers.
The pointwise estimate needed below is
\begin{equation}\label{eq:five-profile-pointwise-certificate}
 q_+(v)\le C\max_i v_i
 -\lambda\sum_{S\in\mathcal S}
       \bigl(\Delta_S(v)-\delta M_S(v)\bigr)-\sigma z(v),
\end{equation}
where
\[
 \lambda=\frac{19}{188}=\frac{3\beta}{4},\quad
 C_0=\frac{96}{47},\quad
 C=C_0-4\lambda\delta=\frac{68768}{34263},\quad
 \sigma=\frac{125}{94}-\frac{304}{34263}>1.
\]
Here is a complete indicator check.
If the root $0$ is present and $\ell$ leaves are present, direct
substitution gives
\[
 q_+(1_J)=A+2\beta\ell-\beta\binom{\ell}{2},\qquad
 \sum_S\Delta_S(1_J)=
 \begin{cases}4,&\ell=0,\\1,&\ell=1,\\0,&\ell\ge2.\end{cases}
\]
If the root is absent and $\ell\ge1$ leaves are present, it gives
$q_+(1_J)=B\ell-\beta\binom{\ell}{2}$ and
$\sum_S\Delta_S(1_J)=0$.  Consequently all nonempty indicator types
are as follows:
\[
\begin{array}{c|c|c|c|c}
0\in J&\ell&
q_+(1_J)+\lambda\sum_S\Delta_S(1_J)&
\sum_SM_S(1_J)&z(1_J)\\ \hline
\text{yes}&0&C_0&4&0\\
\text{yes}&1&C_0-\beta/4&4&0\\
\text{yes}&2&C_0&4&0\\
\text{yes}&3&C_0&4&0\\
\text{yes}&4&C_0-\beta&4&0\\
\text{no}&1&B&3&1\\
\text{no}&2&2B-\beta&4&0\\
\text{no}&3&3B-3\beta&4&0\\
\text{no}&4&C_0&4&0
\end{array}
\]
The leaf-only values increase with $\ell$, since
$B-3\beta=29/94>0$, and the last value is
$4B-6\beta=C_0$.  The singleton-leaf row is an equality in
\eqref{eq:five-profile-pointwise-certificate}, because
$\sigma=C_0-B-\lambda\delta$.  Every other row follows from
$C=C_0-4\lambda\delta$.  The empty indicator has all entries zero.
Integration over layers proves the certificate for every $v\ge0$.

Apply the certificate to
$v_i^+(x)=g_i f(x-h_i)$.  Its four sums of $\Delta_S$ are $D_S$,
its sums of $M_S$ are $P_S$, and its sum of $z$ is $Z$.
For $v_i^-(x)=g_i f(x+h_i)$, translation of $x$ gives
\[
 \sum_x\Phi_{ij}(v^+(x))=\sum_x\Phi_{ij}(v^-(x)).
\]
For example, translating $x$ by $h_i+h_j$ interchanges the two
minima in the definition of $\Phi_{ij}$.  The upper diagonal has
the exact trace
\[
 A+4rB=\frac{77}{47}+\frac{20}{9}\frac{67}{94}
       =\frac{29}{9}=G.
\]
Consequently
$\sum_xq_-(v^-(x))=\sum_xq_+(v^+(x))$.
The tree inequality on $v^-$ and the certificate on $v^+$ yield
\begin{equation}\label{eq:five-profile-defect}
 Q\le C P-\lambda R-\sigma Z.
\end{equation}

Multiply \eqref{eq:five-profile-defect} by $\tau$, multiply
\eqref{eq:five-profile-stability} by $\lambda$, and add.
The $R$ terms cancel:
\[
 Q\le
 \frac{\tau C+\lambda\alpha}{\tau+\lambda}P
 +\frac{\lambda r^2-\tau\sigma}{\tau+\lambda}Z.
\]
The last coefficient is negative, since $\tau>6$, $\sigma>1$,
$\lambda<1$ and $r<1$.  As $Z\ge0$, dropping it and simplifying
gives \eqref{eq:five-profile-bound}.
\end{proof}

The estimate does not give the exact supremum for this profile.
Nor does its retained surplus have a positive lower bound as a
fraction of $P$, uniformly over integer positions.  To see this,
use the geometric-product construction of
Proposition~\ref{prop:geometric-product} with ratio $r=5/9$.
On $\{0,\ldots,L\}^4$ put
$f_L(x)=r^{x_1+\cdots+x_4}$, with $H=\{0,e_1,e_2,e_3,e_4\}$.
Write $s_L=1+r+\cdots+r^L$.  Counting the interior and the disjoint
boundary faces gives
\[
 P_L=s_L^4+4r^{L+1}s_L^3,\quad
 Q_L=s_L^4+4rs_L^3,\quad
 P_{S,L}=s_L^4+3r^{L+1}s_L^3,\quad D_{S,L}=s_L.
\]
In particular,
\[
 \frac{Q_L}{P_L}\longrightarrow\alpha,\qquad
 0\le\frac{R_L}{P_L}
 \le4(s_L^{-3}-\delta)\longrightarrow0.
\]
For each $L$, encoding the four coordinates in an integer base
larger than $2L+3$ preserves these one-coordinate translations and
all displayed counts.  This establishes the same limitation with
integer supports allowed to vary.

\medskip\noindent\textbf{A limit of this combination of estimates.}
It is natural to ask whether different tree probabilities or
coefficients improve the preceding argument.  The answer is no
for the following precise family.  This is a statement about the
family of certificates, rather than the attainable convolution ratio.

Let $p$ be any probability distribution on the spanning trees on
$\{0,1,2,3,4\}$, without a symmetry assumption.  Give each edge the
weight
\[
 w_{ij}=\kappa_{ij}\Pr_{T\sim p}(\{i,j\}\in T),\qquad
 W_p(v)=\sum_{i<j}w_{ij}\Phi_{ij}(v).
\]
Allow real diagonal coefficients $b_i$ satisfying
$\sum_i b_i g_i=G$, and let
$q_{p,b}(v)=\sum_i b_i v_i-W_p(v)$.
Suppose $\lambda,\sigma\ge0$ and $C\in\mathbb R$ give a valid
pointwise certificate
\begin{equation}\label{eq:five-profile-certificate-family}
 q_{p,b}(v)\le C\max_i v_i
 -\lambda\sum_{S\in\mathcal S}
       \bigl(\Delta_S(v)-\delta M_S(v)\bigr)-\sigma z(v)
 \qquad(v\ge0).
\end{equation}
The diagonal normalization and the same reflected pair identities
then give \eqref{eq:five-profile-defect} with these coefficients.

\begin{proposition}\label{prop:five-profile-certificate-barrier}
For every certificate in
\eqref{eq:five-profile-certificate-family},
\[
 \frac{\tau C+\lambda\alpha}{\tau+\lambda}\ge\gamma.
\]
The coefficients in the proof of
Theorem~\ref{thm:five-profile-resolvent} attain equality and make the
remaining $Z$ coefficient negative.  Consequently $\gamma$ is the
best uniform bound obtained by this family when its two estimates
are combined and the remaining $R,Z$ terms are discarded using
only their nonnegativity.
\end{proposition}

\begin{proof}
Put $N=1566459$.  Use the following probability law on nonempty
indicator subsets $J$.  Each named subset has the displayed
probability:
\[
\begin{array}{c|c|c}
J&\text{number of subsets}&\Pr(J)\\ \hline
\{0\}&1&138704/N\\
\{0,i,j\},\ 1\le i<j\le4&6&202220/N\\
\{0,i,j,\ell\},\ 1\le i<j<\ell\le4&4&21715/N\\
\{1,2,3,4\}&1&127575/N
\end{array}
\]
These probabilities are positive and sum to one.  Their root
marginal is $u=1438884/N$ and each leaf marginal is $ru$.
Moreover, every edge has the same expected weighted pair value:
\[
 \mathbb E\bigl[\kappa_{ij}\Phi_{ij}(1_J)\bigr]
 =e=\frac{373225}{N}\qquad(i<j).
\]
For a root--leaf edge, the expected value is
$r\Pr(0,i\in J)$; for a leaf--leaf edge it is
$\Pr(i,j\in J)$.  These give the displayed common value directly.
Every spanning tree has four edges, so
$\mathbb E W_p(1_J)=4e$, for every distribution $p$.
The diagonal normalization gives
$\mathbb E\sum_i b_i1_J(i)=uG$.

Let $a=138704/N$.  Every atom of this law has
\[
 \max_i1_J(i)=1,\qquad
 \sum_SM_S(1_J)=4,\qquad z(1_J)=0.
\]
Only $J=\{0\}$ contributes to $\sum_S\Delta_S(1_J)$, and its
contribution is four.  Averaging
\eqref{eq:five-profile-certificate-family} therefore yields
\[
 C\ge uG-4e+4\lambda(a-\delta).
\]
The chosen rational probabilities satisfy
\[
 uG-4e=\gamma,\qquad
 4\tau(a-\delta)=\gamma-\alpha.
\]
Substitution proves
$C\ge\gamma+\lambda(\gamma-\alpha)/\tau$ and the claimed bound.
The certificate already established attains it.

For completeness, using less of
\eqref{eq:five-profile-stability} cannot improve the conclusion.
If that estimate has weight $t$ in a convex combination, eliminating
$R$ by its nonnegativity requires
$t\tau\le(1-t)\lambda$.  Since the preceding lower bound gives
$C\ge\gamma>\alpha$, the coefficient $(1-t)C+t\alpha$ is smallest
at $t=\lambda/(\tau+\lambda)$.  Requiring the $Z$ coefficient to
be nonpositive can only further restrict the allowed combinations.
\end{proof}

The indicator law is an averaging witness for the stated pointwise
certificate family.  It does not assert that these events occur as
translated levels of one convolution input.  Further relations
between the translates, or estimates outside this family, may
therefore improve the profile bound.  The exact value for this
profile, the general value of $K_5$, and historical priority remain
unresolved here.

\paragraph{A strict bound for finite inputs.}
The order of the integer sites imposes an additional constraint on
actual translated inputs.  It gives the following refinement of
Theorem~\ref{thm:five-profile-resolvent}.

\begin{proposition}\label{prop:five-profile-finite-gap}
Under the assumptions of Theorem~\ref{thm:five-profile-resolvent},
\begin{equation}\label{eq:five-profile-finite-gap}
 \|f\star\widetilde{g1_H}\|_1
 \le \gamma\|f\star g\|_1
       -\frac{1425}{38678}\,t\|f\|_\infty.
\end{equation}
In particular, the ratio in \eqref{eq:five-profile-bound} is strictly
less than $\gamma$ for every nonzero finite input.
\end{proposition}

\begin{proof}
First use the normalized, restricted notation from the proof of
Theorem~\ref{thm:five-profile-resolvent}, so $g_0=1$ and $g_i=r=5/9$
on the four leaves.  For $v\ge0$, let $\ell_1\ge\ell_2$ be its two
largest leaf coordinates, including zeros and ties, and put
\[
 y(v)=\bigl(\min(v_0,\ell_1)-\ell_2\bigr)_+.
\]
Its layer interpretation is the mass of indicators with the root
present and exactly one active leaf.  The existing $z(v)$ is the
mass with the root absent and exactly one active leaf; explicitly,
\[
 z(v)=\bigl(\ell_1-\max(v_0,\ell_2)\bigr)_+,\qquad
 y(v)+z(v)=\ell_1-\ell_2.
\]
With $v_i^+(x)=g_i f(x-h_i)$, set
\[
 Y=\sum_x y(v^+(x)),\qquad T_1(f)=Y+Z.
\]
Thus $T_1(f)$ counts all singleton-leaf layers, regardless of
root membership.

At each threshold $0<s<r\|f\|_\infty$, the finite set
$A_s=\{x:r f(x)>s\}$ is nonempty.  Its four leaf translates are
$A_s+h_i$, $1\le i\le4$.  The least point in these translates,
$\min A_s+\min_i h_i$, belongs only to the translate with the
least shift.  The greatest point, $\max A_s+\max_i h_i$, likewise
belongs only to the translate with the greatest shift.  The two
points are distinct because the leaf shifts are distinct.
Both therefore contribute to a singleton-leaf layer.  Integrating
over $s$ gives
\begin{equation}\label{eq:five-profile-extreme-threshold}
 T_1(f)\ge2r\|f\|_\infty.
\end{equation}
This also holds for $f=0$.  Strict superlevel sets handle ties,
and endpoint thresholds do not affect the integral.

The root-plus-one-leaf row of the indicator table has unused slack
\[
 \eta=\beta-\lambda=\frac{\beta}{4}=\frac{19}{564}.
\]
All the other rows retain nonnegative slack after the existing
$\sigma z(v)$ term.  Since $y(v)$ also has an exact layer integral,
the proof of \eqref{eq:five-profile-pointwise-certificate} gives
the stronger inequality
\[
 q_+(v)\le C\max_i v_i
  -\lambda\sum_{S\in\mathcal S}
       \bigl(\Delta_S(v)-\delta M_S(v)\bigr)
  -\eta y(v)-\sigma z(v).
\]
The same exact-trace reflection argument therefore yields
\[
 Q\le CP-\lambda R-\eta Y-\sigma Z.
\]
Combine this inequality with \eqref{eq:five-profile-stability},
putting weight
\[
 w=\frac{\tau}{\tau+\lambda}=\frac{19035}{19339}
\]
on the former.  The $R$ terms cancel, leaving
\[
 Q\le\gamma P-w\eta Y-
          \bigl(w\sigma-(1-w)r^2\bigr)Z.
\]
The coefficient of $Z$ is greater than $w\eta$, since
$\tau(\sigma-\eta)>\lambda r^2$.  For example, $\sigma>1$,
$\eta<1/9$ and $\tau>6$ make the left side greater than $16/3$,
whereas $\lambda<1/9$ and $r<1$ make the right side less than $1/9$.
Consequently
\begin{equation}\label{eq:five-profile-singleton-gap}
 Q\le\gamma P-w\eta T_1(f)
 \le\gamma P-\frac{1425}{38678}\|f\|_\infty,
\end{equation}
where the last step uses \eqref{eq:five-profile-extreme-threshold}
and
\[
 2rw\eta
 =\frac{10}{9}\frac{19035}{19339}\frac{19}{564}
 =\frac{1425}{38678}.
\]
Restoring the common scale $t>0$ multiplies the restricted numerator
and denominator by $t$.  The original denominator satisfies
$\|f\star g\|_1\ge tP$, so $\gamma>0$ gives
\eqref{eq:five-profile-finite-gap}.
\end{proof}

For the normalized, restricted forward mass $P$, this also gives
\[
 \frac{P}{\|f\|_\infty}\le\frac{115425}{21172}
 \quad\Longrightarrow\quad Q\le2P
 \qquad(f\ne0),
\]
because $\gamma-2=10586/(81\cdot19339)$.  The sufficient condition
is not asserted for every input.

Every atom of the probability law used in the certificate-family
barrier has zero, two, three or four active leaves.  Its
singleton-leaf mass is zero.  By
\eqref{eq:five-profile-extreme-threshold}, that law cannot be the
exact normalized forward layer law of a nonzero finite input.
This constraint does not change the formal averaging argument
that proves the stated certificate-family optimum.

The strict finite-input estimate does not give a smaller uniform
coefficient.  Place widely separated identical copies of any fixed
nonzero finite input so that their forward and reflected output
supports are disjoint.  Both $P$ and $Q$ multiply by the number
of copies, while $\|f\|_\infty$ is unchanged.  Thus
$\|f\|_\infty/P$ has no positive universal lower bound.
The quantity $T_1(f)$ also multiplies under this operation; the
argument therefore concerns the peak-norm consequence, not every
possible use of \eqref{eq:five-profile-singleton-gap}.
It neither shows that the formal witness can be approached nor
determines the true profile supremum.
